\slajdid{10-s003}
\begin{frame}[label=10-s003]
\gusttekst\setlength{\abovedisplayskip}{3pt}\setlength{\belowdisplayskip}{3pt}\setlength{\abovedisplayshortskip}{2pt}\setlength{\belowdisplayshortskip}{2pt}\begin{columns}[c,onlytextwidth]\column{.2\textwidth}Zener dioda\column{.27\textwidth}\centering\begin{tikzpicture}[x=.85cm,y=.65cm,font=\sitneoznake,>=Latex]\ctikzset{bipoles/length=.65cm}\draw(-1,0)--(-.3,0);\draw(-.3,-.3)--(-.3,.3)--(.3,0)--cycle;\draw(.3,0)--(1,0);\draw(.15,-.43)--(.3,-.3)--(.3,.3)--(.45,.43);\node[below]at(-.55,0){A};\node[below]at(.55,0){C};\end{tikzpicture}
\column{.51\textwidth}Direktna polarizacija: kao „obična“ dioda\end{columns}\medskip
\begin{columns}[c,onlytextwidth]\column{.2\textwidth}\column{.27\textwidth}\centering\begin{tikzpicture}[x=.85cm,y=.65cm,font=\sitneoznake,>=Latex]\ctikzset{bipoles/length=.65cm}\draw(-1,0)node[ocirc]{}node[left]{A}--(-.08,0);\draw(.08,0)--(1,0)node[ocirc]{}node[right]{C};\draw(-.08,-.2)--(-.08,.2);\draw(.08,-.5)--(.08,.5);\node[above]at(0,.65){$V_Z$};\node at(.3,.35){$+$};\draw[<-](.4,.3)--(.85,.3)node[midway,above]{$I_Z$};\end{tikzpicture}
\column{.51\textwidth}Inverzna polarizacija:\par ostvareni uslovi\par Zenerov proboj\par model naponskog izvora\end{columns}\bigskip
\begin{columns}[c,onlytextwidth]\column{.2\textwidth}Šotki dioda\column{.27\textwidth}\centering\begin{tikzpicture}[x=.85cm,y=.65cm,font=\sitneoznake,>=Latex]\ctikzset{bipoles/length=.65cm}\draw(-1,0)--(-.3,0);\draw(-.3,-.3)--(-.3,.3)--(.3,0)--cycle;\draw(.3,0)--(1,0);\draw(.15,-.3)--(.3,-.3)--(.3,.3)--(.45,.3);\node[below]at(-.55,0){A};\node[below]at(.55,0){C};\end{tikzpicture}
\column{.51\textwidth}\end{columns}\bigskip
Šotki diodu obrazuje spoj metal poluprovodnik i zbog toga je napon direktne polarizacije kod ovakvih dioda manji. Za Šotki diode važe isti modeli kao i za pn diode, osim što ćemo u daljoj analizi smatrati da je pri direktnoj polarizaciji Šotki diode $V_{AC}=V_S=0.3V$.
\par\medskip Znatno su „brže“ od običnih dioda. Nema nagomilanih nosilaca.
\end{frame}
